\paragraph{Prefix extension as an off-diagonal collision.}
For \(f:W\to Y\) and \(t\in Y\), take two tagged copies of the candidate
space and define
\[
 G(0,w)=(w,t),\qquad G(1,w)=(w,f(w)).
\]
The first coordinate forces colliding candidates to be identical. Two
distinct tagged points with that candidate collide exactly when
\(f(w)=t\). Thus the ordered off-diagonal of \(G\) is two copies of
\(f^{-1}(t)\), one per choice of the first tag. In particular a unique
witness already contributes two ordered collisions. For any binary
verifier take \(f(w)=(V(x,w),w_{1:k})\) and \(t=(1,u)\). The
off-diagonal is inhabited exactly when \(u\) extends. Lean proves the
equivalence. For an affine verifier \(Qw=c\), use instead the linear
map \(w\mapsto(Qw,w_{1:k})\) and target \((c,u)\). Its fiber consists
of the original extending witnesses.

There is also an exact binary affine system for \emph{any} verifier
when exponential size is allowed. After fixing \(u\), introduce one
variable \(\lambda_s\) per remaining suffix and impose
\[
 \sum_{s\in\{0,1\}^{n-k}}\lambda_s=1,\qquad
 (1-V(x,us))\lambda_s=0\quad\text{for each }s .
\]
An extending suffix gives its unit vector. Conversely, at least one
coordinate of any solution is nonzero, and every such coordinate
indexes an extending suffix. This is an affine fiber, so the same
tagged graph construction realizes its feasibility as a collision.
It uses \(2^{n-k}\) variables. The following family admits a
polynomial-size construction.

\begin{proposition}[An affine consistency test for every 2-CNF prefix]
\label{prop:two-cnf-prefix}
Let \(F\) be a conjunction of clauses of size at most two on \(n\)
explicitly ordered variables. From \(F\) and any \(k\)-bit prefix \(u\),
\(0\le k\le n\), one can construct in polynomial time a binary affine
system on \(n\) variables that is consistent exactly when \(u\)
extends to a satisfying assignment of \(F\).
\end{proposition}
\begin{proof}
Write \(\tau(x_i)=v_i\) and \(\tau(\neg x_i)=1+v_i\) over \(\Ftwo\).
Use the implication graph with edges \(\neg a\to b,\ \neg b\to a\)
for each \(a\lor b\), repeating a literal for a unit clause. Include
\(0=1\) if there is an empty clause. In each strongly connected
component choose a representative \(r\) and impose
\(\tau(\ell)=\tau(r)\) for its other literals. These rows depend only
on \(F\). From the literals set true by the prefix, take their
forward reachability closure \(T\), including the starting literals.
Append \(\tau(\ell)=1\) for each \(\ell\in T\).

Every extending assignment satisfies the rows. Conversely, consistency
rules out a component containing a complementary pair, since its two
truth values would be equal. It also gives \(T\cap\neg T=\varnothing\).
Assign \(T\) true and its complements false. If a clause contains a
literal assigned false, the edge from its true complement forces the
other literal into \(T\). Only clauses wholly on unassigned variables
still require an assignment. Their implication graph is a subgraph of
the original graph, so no component contains a complementary pair.
The strong-component criterion for 2-SAT gives a satisfying completion
\cite[Theorem~1]{aspvallplassTarjan}.

Reachability from the \(2n\) literals in the base graph can be
precomputed once by graph search. Each prefix uses at most \(2n\)
component rows, \(2n\) pins and one empty-clause row. The resulting
dense binary system has polynomial size and can be tested by elimination.
\end{proof}
For \(F=x_1\lor x_2\), prefix \(0\) forces \(\neg x_1\) and \(x_2\).
The affine system is \(v_1=0,\ v_2=1\). Prefix \(00\) gives incompatible
pins, and prefix search returns \(01\) from the empty prefix. The
empty-prefix system accepts the auxiliary point \(00\) too. Its
off-diagonal detects the existence of an extension, while the
original verifier's off-diagonal counts the actual extensions.
The three satisfying full prefixes of \(x_1\lor x_2\) cannot be the
projection of a single binary affine solution set, since every
nonempty affine image has power-of-two cardinality. Our construction
uses the prefix to determine which rows to append.
